%%%%%%%%%%%%%%%%%%%%%%%%%%%%%%%%%%%%%%%%%%%%%%%%%%%%%%%%%%%%%%%%%%%%%%
% Lớp 10 | Đại số tổ hợp | Quy tắc cộng-quy tắc nhân
% Dạng 0D8?1-4: Bài toán dùng quy tắc bù trừ
% Số câu: 5  (sinh tự động bởi tach_cau.py; từ nay thầy sửa trực tiếp file này)
%%%%%%%%%%%%%%%%%%%%%%%%%%%%%%%%%%%%%%%%%%%%%%%%%%%%%%%%%%%%%%%%%%%%%%

%%%%%%%%%%%%%%%% TRẮC NGHIỆM 4 PHƯƠNG ÁN %%%%%%%%%%%%%%%%

% nguồn: MaHoaNH_TracNghiem_Lop10
\begin{ex}%[0D8N1-4]%[ma:4cfeefaf]
	Có $10$ cặp vợ chồng đi  dự tiệc. Tính số cách chọn một người đàn ông và một người đàn bà trong bữa tiệc phát biểu ý kiến sao cho hai người đó không là vợ chồng.
	 \choice
		 {$19$}
		{\True $90$} 
		 {$100$}
		 {$20$}
	\loigiai{ 
		Chọn $1$ người đàn ông phát biểu có $10$ cách.\\ 
		Chọn $1$ người đàn bà phát biểu  có $10$ cách.\\
		Số cách chọn một người đàn ông và một người đàn bà trong bữa tiệc phát  biểu ý kiến sao cho hai người đó không là vợ chồng là $10\cdot10-10=90$ cách.
	 }
\end{ex}

% nguồn: MaHoaNH_TracNghiem_Lop10
\begin{ex}%[0D8V1-4]%[ma:1d83edfa]
	Từ các chữ số $0$, $2$, $3$, $5$, $6$, $8$ có thể lập được bao nhiêu số  tự nhiên gồm $6$ chữ số đôi một khác nhau trong đó hai chữ số $0$ và $5$ không đứng cạnh nhau.
	 \choice
		 {\True $384$}
		 {$120$}
		 {$216$}
		{$600$} 
	\loigiai{ 
		Số các số có $6$ chữ số  đôi một khác nhau lập từ $0$, $2$, $3$, $5$, $6$, $8$ là $5\cdot5\cdot4\cdot3\cdot2\cdot1=600$ số.\\
		Ta đi  đếm số các số có chữ số $0$ và $5$ đứng cạnh nhau.\\
		Giả sử  số cần tìm có dạng $\overline{abcdef}$.\\
		 Khi đó, bộ $05$ có thể thay thế các cặp chữ $bc$, $cd$, $de$, $ef$. Ta có $4$ cách.\\
		Ứng với mỗi cách thay thế, ta có $4\cdot3\cdot2\cdot1$ cách sắp xếp các chữ số $2$, $3$, $6$, $8$ vào các chữ còn lại. Do đó, số các số lập  được là $4\cdot4\cdot 3\cdot2\cdot1=96$ số.\\
		Trường hợp ta xếp bộ $50$, ta có thể thay thế cho các cặp chữ $ab$, $bc$, $cd$,  $de$, $ef$. Ta có $5$ cách. Tương tự cách xếp bộ $50$, số các số lập được là $5\cdot4\cdot 3\cdot2\cdot1=120$ số.\\
		Vậy số các số có $6$ chữ số  khác nhau thỏa $0$ và $5$ đứng cạnh nhau là $120+96=216$ số.\\
		Theo nguyên lý bù trừ, số các  số có $6$ chữ số khác nhau thỏa $0$ và $5$ không đứng cạnh nhau là $600-216=384$ số.
	 }
\end{ex}


%%%%%%%%%%%%%%%% ĐÚNG/SAI %%%%%%%%%%%%%%%%

% nguồn: MaHoaNH_DungSai_Lop10
\begin{ex}%[0D8H1-4]%[ma:8ae838f2]
	Một lớp học có $36$ học sinh chia thành $3$ tổ. Tổ $1$ có $11$ học sinh, tổ $2$ có $12$ học sinh, tổ $3$ có $13$ học  sinh. Cần chọn $2$ học sinh đại diện lớp đi dự hội nghị.
	\choiceTF 
		{\True Có $132$ cách  chọn hai học sinh gồm một học sinh ở tổ $1$ và một học sinh ở tổ $2$}
		{\True Có $143$ cách chọn hai học sinh gồm  một học sinh ở tổ $1$ và một học sinh ở tổ $3$}
		{\True Có $156$ cách chọn hai học sinh gồm $1$ học sinh ở  tổ $2$ và một học sinh ở tổ $3$}
		{Có $362$ cách  chọn hai học sinh cùng tổ}
	\loigiai{ 
		Chọn  hai học sinh bất kỳ có $36\cdot 35:2=630$ cách.\\
		Có $11\cdot 12=132$ cách chọn hai học sinh gồm  một học sinh ở tổ $1$ và một học sinh ở tổ $2$.\\
		Có $11\cdot 13=143$ cách chọn hai học sinh gồm một học sinh ở tổ $1$  và một học sinh ở tổ $3$.\\
		Có $12\cdot 13=156$ cách chọn hai học sinh gồm một học sinh ở tổ $1$ và một học sinh  ở tổ $3$.\\
		Vậy có  $630-132-143-156=199$ cách chọn hai học sinh cùng một tổ.
	 }
\end{ex}

% nguồn: MaHoaNH_DungSai_Lop10
\begin{ex}%[0D8H1-4]%[ma:f1b190e1]
	Một học sinh có $6$ cây viết đỏ khác nhau, $8$ cây viết xanh khác nhau, $7$ cây  viết đen khác nhau. Học sinh chọn hai cây viết cùng màu để tặng.
	 \choiceTF
		{\True Có  $48$ cách chọn hai cây viết gồm $1$ đỏ và $1$ xanh}
		{\True Có $42$ cách chọn  hai cây viết gồm $1$ đỏ và $1$ đen}
		{\True Có $56$ cách chọn hai cây viết gồm $1$ xanh và $1$ đen} 
		{Có $94$ cách chọn  hai cây viết cùng màu}
	 \loigiai{
		Chọn hai cây viết bất kỳ có  $21\cdot 20:2=210$ cách.\\
		Chọn hai cây viết khác màu có màu đỏ xanh ($1$ đỏ, $1$  xanh) có $6\cdot 8=48$ cách.\\
		Chọn hai cây viết  khác màu có màu đỏ đen ($1$ đỏ, $1$ đen) có $6\cdot 7=42$ cách.\\
		Chọn hai cây viết khác màu có màu xanh đen ($1$ xanh, $1$ đen) có $8\cdot  7=56$ cách.\\
		Chọn hai cây viết khác màu có $6\cdot 8+6\cdot  7+7\cdot 8=146$ cách.\\
		 Vậy có $210-48-42-56=64$ cách chọn hai cây viết cùng màu.
	 }
\end{ex}


%%%%%%%%%%%%%%%% TRẢ LỜI NGẮN %%%%%%%%%%%%%%%%

% nguồn: MaHoaNH_TraLoiNgan_Lop10
\begin{bt}%[0D8N1-4]%[ma:7e41415b]
	Từ $6$ chữ số $1$,  $2$, $3$, $4$, $5$, $6$ có thể lập được bao nhiêu số các số gồm $3$ chữ số đôi một khác nhau không chia hết cho $5$?
	\shortans{100} 
	 \loigiai{
		Gọi số có dạng $\overline{a b c}$ là số có ba chữ số đôi một khác nhau và chia hết cho $5$ với  $a$, $b$, $c\in\{1; 2; 3; 4; 5; 6\}$ thì số cần tìm có dạng $ \overline{ab5}$.\\
		Chọn $c$ có $1$ cách chọn  ($c=5$);\\
		Khi đó chọn  $a$ có $5$ cách chọn (khác $d$), chọn $b$ có $4$ cách chọn.\\
		Vậy có tất cả $1\cdot5\cdot4=20$ số chia hết cho $5$  được lấy từ $\{1; 2; 3; 4; 5; 6\}$.\\
		Có $6\cdot 5\cdot 4=120$ số có ba  chữ số khác nhau được lấy từ $\{1; 2; 3; 4; 5; 6\}$.\\
		 Vậy có $120-20=100$ số cần tìm.
	} 
\end{bt}
